Read in Gauge Theory Gravity, the metric below is the derived metric of the gravitational field, its variation is the variation of $\hb$ with $\Omega$ torsion-free, and the Ricci and Riemann tensors are contractions of the field strength (Fact \claimref{synthesis.imported-gtg-dictionary}, clauses (i) and (iv)). \emph{Conventions.} Signature $(-,+,+,+)$, Ricci tensor positive on spheres, as in the check.

\emph{(i).} Write $f = \sum_{k \geq 2} f_{k}$ with $f_{k}$ homogeneous of degree $k$ in $R_{ab}$, its coefficients built from the metric. The field equations of $\sqrt{-g}\,f$ are obtained by varying $g^{ab}$: the variation of $\sqrt{-g}$ and of the metric contractions gives terms proportional to $f_{k}$ or to $\partial f_{k}/\partial g^{ab}$ at fixed $R_{ab}$, both of degree $k \geq 2$ in the Ricci tensor; the variation of $R_{ab}$ gives $P^{ab}\,\delta R_{ab}$ with $P^{ab} = \partial f/\partial R_{ab}$ of degree at least one in the Ricci tensor, and since $\delta R_{ab}$ is a combination of second covariant derivatives of $\delta g$, integration by parts leaves second covariant derivatives of $P^{ab}$ and products of $P^{ab}$ with curvature. On a Ricci-flat metric $P^{ab}$ vanishes identically, so all its derivatives vanish, and every term of the field equations of $f$ vanishes; the Einstein part vanishes because the metric is Ricci-flat. The check computes the full curvature of $g = \operatorname{diag}(-h, 1/h, r^{2}/(1 - x^{2}), r^{2}(1 - x^{2}))$, $x = \cos\theta$, from the Christoffel symbols (exact rational functions of $r$, Taylor series in $x$ about $\tfrac12$ to the order the second derivatives need): for $h = 1 - 2M/r$ it finds $R_{ab} = 0$ and $K = 48M^{2}/r^{6}$; for flat space the whole Riemann tensor vanishes. The surface gravity is $\tfrac12 h'(2M) = 1/4M$. Since $\sqrt{K/12} = 2M/r^{3}$, the surface of (C1), $2M/r^{3} = 1$ in Planck units, is a level set of a curvature invariant; in a region where $K < 12$, in particular in flat space, it is absent.

\emph{(ii).} By (W) and clause (b) of Fact \claimref{synthesis.imported-wald-noether-charge-entropy}, $S = -2\pi\oint E^{abcd}\epsilon_{ab}\epsilon_{cd}\,dA$ with $E^{abcd} = \partial L/\partial R_{abcd}$ and $\epsilon_{ab}$ the binormal. For the Einstein term this gives $A/4G$. For $f$, $\partial f/\partial R_{abcd}$ is $P^{bd}g^{ac}$ symmetrised with the symmetries of the Riemann tensor, with $P^{bd} = \partial f/\partial R_{bd}$ of degree at least one in the Ricci tensor; it vanishes on the Ricci-flat horizon, so $S = A/4G$ for every $f$.

\emph{(iii).} Theorem \claimref{singularity.rate-discharge-hypothesis-iii} derives the temperature $1/8\pi M$ from (P), (R) and the geometry of the Euclidean section, whose period $2\pi/\kappa$ depends only on the metric, so it is the same under every $L_{f}$ by (i). Theorem \claimref{scrambling.mode-count} counts modes on the surface of (C1) with the density of (C2) over the slice volume of the interior metric: again geometry only. Proposition \claimref{synthesis.quarter-from-temperature} integrates $dS = dM/T$ along the Schwarzschild family to $4\pi M^{2} = A/4$, which by (ii) is also the Wald entropy under every $L_{f}$. So any derivation that uses only these inputs reaches the same conclusions for every $f$. The field equations differ: for $f = \beta R_{ab}R^{ab}$ the term $P^{ab}\delta R_{ab}$ contributes $\beta\,\Box R_{ab}$ and further derivative terms, so with matter the equations are of fourth order and are not $G_{ab} = 8\pi G T_{ab}$; a derivation of Einstein's equations from data shared by all $L_{f}$ would derive them in every $L_{f}$, which is false for $f \neq 0$.

\emph{(iv).} The check computes the Ricci tensor of the Reissner-Nordstr\"om metric $h = 1 - 2M/r + Q^{2}/r^{2}$: $R^{t}{}_{t} = R^{r}{}_{r} = -Q^{2}/r^{4}$, $R^{x}{}_{x} = R^{\phi}{}_{\phi} = Q^{2}/r^{4}$, off-diagonal components zero, nonzero at the horizon $r_{+}$ (at $M = \tfrac32$, $Q^{2} = \tfrac54$, $r_{+} = \tfrac52$). For $f = \beta R_{ab}R^{ab}$, $E^{abcd}\epsilon_{ab}\epsilon_{cd} = 2\beta R^{bd}\epsilon_{ab}\epsilon^{a}{}_{d} = -2\beta\,(R^{t}{}_{t} + R^{r}{}_{r})$ with the binormal $\epsilon_{tr} = -\epsilon_{rt} = 1$ (here $-g_{tt}g_{rr} = 1$), normalised by $\epsilon_{ab}\epsilon^{ab} = -2$ (the check verifies this contraction as a pointwise identity, outside the horizon where these coordinates are regular), so the Wald entropy gains $4\pi\beta\,(R^{t}{}_{t} + R^{r}{}_{r})A$, which is $-8\pi\beta Q^{2}A/r_{+}^{4} \neq 0$ there. In Jacobson's derivation the heat flux through a local Rindler horizon is $\int T_{ab}k^{a}k^{b}$ and, through Einstein's equations, equals the focusing term $R_{ab}k^{a}k^{b}$; the Clausius relation is imposed for every point and every null $k^{a}$, which is where the Ricci tensor is nonzero and where (by the computation just made) the Ricci-dependent part of a horizon entropy is visible. For the remark: (R) is a premise about stationary Killing horizons and (C1) about Planck-curvature surfaces; by (i) a black hole's (C1) surface is absent from any region with $\sqrt{K/12} < 1$. The remark is a reading of the store, not part of the proved clauses.
